\scanpage{268}


Pour le cas ``dégénéré'', notons

\[
\begin{aligned}
\sigma\in T_H&\Longrightarrow U_0=U_1=\{1\},\quad \deleted{\widetilde Z}\ T_H\ \text{un }\deleted{\text{sous-}}\text{groupe}\\
\Updownarrow&\hspace{15em}\text{commutatif avec}\\
\deleted{\unsure{\exists\lambda\in Z,\ t_0=\lambda s_0}}&\hspace{15em}t\in T_H\longrightarrow t^2=1\\
t_0\in Z.s_0&\qquad\text{enfin }\widetilde Z=\operatorname{Cent}_H(T_H).\\
\text{i.e. }\underset{D_0}{\underbrace{t_0^{\unsure{\leftrightarrows|\leftrightarrow|\Diamond}}}}
=\underset{D_1}{\underbrace{s_0^{\leftrightarrows}}}&
\end{aligned}\tag*{(155)}
\]

% slanted
\margin{\unsure{implique|impliques} donc que $S$ est un \uline{torseur} sous $H$ ($\Longleftrightarrow U_0=\{1\}\Longrightarrow U_1=\{1\}$),}

\deleted{\unsure{ainsi} Notons que la} Plus précisément, la donnée d'un multigraphe à homothéties et groupe simplement transitif d'autom \inserted{et \deleted{point} épinglé par $(\unsure{a|s|s_0},b)\in A_0$,} \deleted{tel que $\sigma\in T_H$ i.e} et tel que $s_0^{\leftrightarrows}=t_0^{\leftrightarrows}$ i.e $\sigma\in T_H$, \uline{équivaut} à la donnée de $H$, du sous-groupe commutatif $T_H$ tel que $T_H^2=\{1\}$, et de $\sigma\in T_H$ -- et on aura alors $\widetilde Z=\operatorname{Cent}_H(T_H)$.

\deleted{Considérons}\inserted{Notons} le cas, encore plus dégénéré

\[
\sigma=1\ .\Longleftrightarrow s_0=t_0\Longleftrightarrow
\begin{aligned}& A_0\ \text{est formé des parties}\\[-2pt]&\text{ponctuelles de }S.\end{aligned}\tag*{(156)}
\]

\begin{center}\rule{2.1cm}{0.3pt}\end{center}

\unsure{1} Digression dans le cas d'un multigraphe avec \inserted{gr de} homothéties \inserted{$Z$}, et groupe \deleted{d'}\inserted{de $Z$-}autom. $H$ simplement transitif sur \deleted{les arêtes} un des $\overrightarrow A_\lambda$ (donc sur tous) -- Fixons nous une origine $s_0$ et un rayon incident \unsure{à|:} $s_0$

\[
\deleted{\unsure{s_0,\ \text{donc}|\text{donc }s_0}}\quad s_0\in S,\ D_1\in S^{\leftrightarrows},\quad D_0=Z.s_0\ \text{et }D_1\ \text{incidents dans }S^{\leftrightarrows}\ ;\tag*{(157)}
\]

donc on a

\[
\begin{aligned}
S&\simeq H/U_0\qquad\text{où}\quad U=H_{s_0}\\
&\text{(isom de }H\text{-ensembles)}\ ,
\end{aligned}\tag*{(158)}
\]

\deleted{fixons nous en plus un rayon} de plus $T_G$, $T_0$ \deleted{\unsure{G_{\leftrightarrows}} et} $G_S$ sont bien déterminés par (137), \deleted{et $N_1$} ou (90) (\unsure{93|92|97}), donc aussi $T_H=T_G\cap H$ :

\[
\deleted{T_H\subset H,\ T_0\subset T_G\in G,}
\]
% archive sheet number: 267
